\subsection{H5: 4-Node Motif FFL-Derivativeness}\label{sec:results_h5}

Analysis of 179 pruned attribution graphs (pruned at the 99th percentile to ensure computational tractability, with at most 700 nodes per graph) identifies 4 universal 4-node motif types (IDs 77, 80, 82, 83). These are the 4-node directed subgraph patterns that appear across all or nearly all capability domains. Functional characterization reveals a striking pattern of FFL derivativeness:

\begin{itemize}
\item \textbf{100\% FFL containment:} Every single instance of every universal 4-node motif contains at least one embedded 3-node FFL subgraph. This is a remarkably clean result: higher-order motif universality is \emph{entirely} explained by FFL composition. No universal 4-node pattern exists independently of FFL structure.
\item \textbf{100\% strict layer ordering:} All 4-node motif instances respect the DAG layer structure of the attribution graph, with nodes spanning increasing layer indices. This is consistent with the feedforward nature of transformer computation and FFL layer ordering.
\item \textbf{Enhanced cross-domain association:} Cram\'{e}r's $V$ values of 0.20--0.33 for 4-node motif-domain associations exceed the 3-node FFL baseline of $V=0.13$. This suggests that 4-node motifs, while FFL-derivative, capture finer-grained structural variation that differentiates domains more strongly than FFLs alone.
\item \textbf{Extended layer span:} 4-node motif instances span significantly more transformer layers than random connected 4-node subgraphs ($p<0.001$), indicating that universal motifs preferentially involve long-range hierarchical connections across the transformer stack.
\end{itemize}

A total of 6,574 motif--graph characterization records were generated across all graphs. Random baselines (connected 4-node subgraphs sampled without regard to motif structure) show zero motif matches in every graph tested, confirming that the identified patterns are genuine structural motifs rather than sampling artifacts. The four universal motif types correspond to distinct FFL extensions: motif 77 (4 edges, fan-out from FFL source), motif 80 (5 edges, chain extension), motif 82 (5 edges, convergent extension), and motif 83 (6 edges, fully connected FFL extension).

This compositional hierarchy resonates with findings from biological network analysis, where higher-order motifs in gene regulatory and neuronal networks are frequently composed of simpler building blocks \cite{milo2004superfamilies, benson2016higher}. The 100\% FFL containment result is stronger than typical biological findings, reflecting the more constrained structure of layered DAGs compared to general directed networks. The enhanced cross-domain association of 4-node motifs (Cram\'{e}r's $V=0.20$--$0.33$ vs.\ $V=0.13$ for FFLs) suggests a practical tradeoff: while FFLs are universal, their 4-node extensions capture finer-grained domain-specific structure (Table~\ref{table:ablation}). Future work could exploit this hierarchy for multi-resolution circuit analysis, using FFLs for broad characterization and 4-node motifs for domain discrimination.

\textbf{Verdict:} H5 is \textbf{supported}. All universal 4-node motifs are fully FFL-derivative with 100\% containment, establishing the 3-node FFL as the fundamental building block from which higher-order motif structure emerges in attribution circuits.